\section{PDDL-Based Crate Navigation}
\label{sec:pddl}

Some maps in the Deliveroo.js environment contain movable crates that can block navigable paths. An agent may push a crate one tile in any cardinal direction, subject to the constraints described in Section~\ref{sec:planning}. We model crate manipulation using PDDL (Planning Domain Definition Language), invoked as a fallback when A* alone cannot find a route.

\subsection{Domain Overview}

The domain \texttt{deliveroo-crates} is specified in STRIPS with typing (\texttt{:strips :typing}) over two object types: \texttt{tile} and \texttt{crate}. The predicate set captures the complete state required for crate-aware navigation:

\begin{itemize}[noitemsep]
  \item \texttt{(at ?t)} --- agent occupies tile $t$;
  \item \texttt{(adj-up/down/left/right ?from ?to)} --- directed adjacency between tiles;
  \item \texttt{(crate-at ?c ?t)} --- crate $c$ is positioned on tile $t$;
  \item \texttt{(crate-free ?t)} --- tile $t$ contains no crate (available for entry or push landing);
  \item \texttt{(crate-space ?t)} --- tile $t$ is structurally valid for crate occupation (not a wall or conveyor).
\end{itemize}

Eight actions are defined. Four \emph{move} actions allow the agent to step into an adjacent crate-free tile. Four \emph{push} actions allow the agent to displace a crate one tile in the push direction, with preconditions: (i) the agent is adjacent to the crate from the opposite side; (ii) the destination is a valid crate space and is currently crate-free. The push effect simultaneously relocates the agent onto the crate's former tile and the crate to its new tile, updating \texttt{crate-free} facts accordingly.

\subsection{Planning With PDDL}

PDDL planning is triggered exclusively when A* fails and the belief system confirms that crates obstruct the route (Section~\ref{sec:planning}). The problem instance is generated by \texttt{buildProblem}: it emits adjacency facts for every non-wall tile pair, \texttt{crate-at} facts for each believed crate, \texttt{crate-free} facts for all unoccupied tiles, and \texttt{crate-space} facts for all structurally valid positions. The goal is a single ground atom \texttt{(at target)}.

Two invocation modes exist. In the \emph{stitched} mode, the planner receives only the minimal entry--exit cluster segment identified by \texttt{findCrateSegment} (Section~\ref{sec:belief}); the A* prefix to the cluster entry and the A* suffix from the cluster exit to the final goal are computed separately and concatenated. In the \emph{whole-trip} fallback mode, the full route from the agent's current position to the goal is submitted when the segmentation heuristic fails (e.g.\ multi-cluster corridors).

Solver selection is runtime-configurable: when the environment variable \texttt{PAAS\_HOST} is set, problems are submitted to the online \texttt{@unitn-asa/pddl-client} solver; otherwise the local binary \texttt{dual-bfws-ffparser} (accessed via \texttt{planutils}) is invoked in a per-call temporary directory to avoid file-system races in multi-agent deployments. Both \texttt{move-*} and \texttt{push-*} steps in the returned plan are translated into a \texttt{move} plan step (the agent's tile after the action), and appended to the A* prefix.

\subsection{A Smart Solution --- Problem Restatement \& Caching}

Since the PDDL planner introduced overhead that was not acceptable for the reactive nature of the environment, we implemented a problem restatement to enable plan caching.

Instead of invoking the planner for a global solution from the agent's current position to the final goal, the plan is decomposed into three steps:
\begin{itemize}[noitemsep]
    \item Move from the current position to the cluster entry tile (A*);
    \item Navigate through the crate cluster (PDDL);
    \item Move from the cluster exit tile to the final goal (A*).
\end{itemize}
The PDDL planner is therefore restricted to the second step alone. Computed maneuvers are stored in a FIFO cache of size 32, keyed by the tuple $(\text{entry},\, \text{exit},\, \{\text{crate positions}\})$. The cache key is independent of the agent's absolute position: it depends only on the relative crate arrangement and the cluster's entry and exit tiles. Consequently, the exact same maneuver plan is reused whenever the agent encounters a structurally identical sub-problem --- including cases where the agent approaches the same cluster from a different direction --- drastically reducing planner invocations and keeping the agent reactive.
